\subsection{Notations and elementary facts}

We will use the following notations and conventions:

\begin{itemize}
\item The symbol $\mathbb{N}$ will denote the set $\left\{  0,1,2,3,\ldots
\right\}  $ of nonnegative integers.

\item The size (i.e., cardinality) of a set $A$ will be denoted by $\left\vert
A\right\vert $.

\item The symbol \textquotedblleft\#\textquotedblright\ means
\textquotedblleft number\textquotedblright. For example, the size $\left\vert
A\right\vert $ of a set $A$ is the \# of elements of $A$.
\end{itemize}

We will need some basics from enumerative combinatorics (see, e.g.,
\cite[\S 8.1]{Newste19} for details, and \cite[Chapters 1 and 2]{19fco} for
more details):

\begin{itemize}
\item \textbf{addition principle = sum rule:} If $A$ and $B$ are two disjoint
sets, then $\left\vert A\cup B\right\vert =\left\vert A\right\vert +\left\vert
B\right\vert $.

\item \textbf{multiplication principle = product rule:} If $A$ and $B$ are any
two sets, then $\left\vert A\times B\right\vert =\left\vert A\right\vert
\cdot\left\vert B\right\vert $.

\item \textbf{bijection principle:} There is a bijection (= bijective map =
invertible map = one-to-one correspondence) between two sets $X$ and $Y$ if
and only if $\left\vert X\right\vert =\left\vert Y\right\vert $.

\item A set with $n$ elements has $2^{n}$ subsets, and has $\dbinom{n}{k}$
size-$k$ subsets for any $k\in\mathbb{R}$.

\item A set with $n$ elements has $n!$ permutations (= bijective maps from
this set to itself).

\item \textbf{dependent product rule:} Consider a situation in which you have
to make $n$ choices (sequentially). Assume that you have $a_{1}$ options
available in choice $1$, and then (after making choice $1$) you have $a_{2}$
options available in choice $2$ (no matter which option you chose in choice
$1$), and then (after both choices $1$ and $2$) you have $a_{3}$ options
available in choice $3$ (no matter which options you chose in choices $1$ and
$2$), and so on. Then, the total \# of ways to make all $n$ choices is
$a_{1}a_{2}\cdots a_{n}$. (This is formalized in \cite[Theorem 8.1.19]%
{Newste19}.)
\end{itemize}

A few words about binomial coefficients are in order:

\begin{definition}
\label{def.binom.binom}For any numbers $n$ and $k$, we set%
\begin{equation}
\dbinom{n}{k}=%
\begin{cases}
\dfrac{n\left(  n-1\right)  \left(  n-2\right)  \cdots\left(  n-k+1\right)
}{k!}, & \text{if }k\in\mathbb{N};\\
0, & \text{else.}%
\end{cases}
\label{eq.def.binom.binom.eq}%
\end{equation}
Note that \textquotedblleft numbers\textquotedblright\ is to be understood
fairly liberally here. In particular, $n$ can be any integer, rational, real
or complex number (or, more generally, any element in a $\mathbb{Q}$-algebra),
whereas $k$ can be anything (although the only nonzero values of $\dbinom
{n}{k}$ will be achieved for $k\in\mathbb{N}$, by the above definition).
\end{definition}

\begin{example}
\label{exa.binom.-1choosek}For any $k\in\mathbb{N}$, we have
\begin{align*}
\dbinom{-1}{k}  &  =\dfrac{\left(  -1\right)  \left(  -1-1\right)  \left(
-1-2\right)  \cdots\left(  -1-k+1\right)  }{k!}\\
&  =\dfrac{\left(  -1\right)  \left(  -2\right)  \left(  -3\right)
\cdots\left(  -k\right)  }{k!}=\dfrac{\left(  -1\right)  ^{k}k!}{k!}=\left(
-1\right)  ^{k}.
\end{align*}

\end{example}

If $n,k\in\mathbb{N}$ and $n\geq k$, then
\begin{equation}
\dbinom{n}{k}=\dfrac{n!}{k!\left(  n-k\right)  !}. \label{eq.binom.fac-form}%
\end{equation}
But this formula only applies to the case when $n,k\in\mathbb{N}$ and $n\geq
k$. Our above definition is more general than it. The combinatorial meaning of
the binomial coefficient $\dbinom{n}{k}$ (as the \# of $k$-element subsets of
a given $n$-element set) also cannot be used for negative or non-integer
values of $n$.

\begin{example}
\label{exa.binom.2n-choose-n}Let $n\in\mathbb{N}$. Then, $\dbinom{2n}%
{n}=\dfrac{1\cdot3\cdot5\cdot\cdots\cdot\left(  2n-1\right)  }{n!}\cdot2^{n}$.
\end{example}

\begin{proof}
[Proof of Example \ref{exa.binom.2n-choose-n}.]We have%
\begin{align*}
\left(  2n\right)  !  &  =1\cdot2\cdot\cdots\cdot\left(  2n\right) \\
&  =\left(  1\cdot3\cdot5\cdot\cdots\cdot\left(  2n-1\right)  \right)
\cdot\underbrace{\left(  2\cdot4\cdot6\cdot\cdots\cdot\left(  2n\right)
\right)  }_{\substack{=\left(  2\cdot1\right)  \cdot\left(  2\cdot2\right)
\cdot\cdots\cdot\left(  2\cdot n\right)  \\=2^{n}\left(  1\cdot2\cdot
\cdots\cdot n\right)  }}\\
&  =\left(  1\cdot3\cdot5\cdot\cdots\cdot\left(  2n-1\right)  \right)
\cdot2^{n}\underbrace{\left(  1\cdot2\cdot\cdots\cdot n\right)  }_{=n!}\\
&  =\left(  1\cdot3\cdot5\cdot\cdots\cdot\left(  2n-1\right)  \right)
\cdot2^{n}n!.
\end{align*}
Now, (\ref{eq.binom.fac-form}) yields%
\begin{align*}
\dbinom{2n}{n}  &  =\dfrac{\left(  2n\right)  !}{n!\left(  2n-n\right)
!}=\dfrac{\left(  2n\right)  !}{n!n!}=\dfrac{\left(  1\cdot3\cdot5\cdot
\cdots\cdot\left(  2n-1\right)  \right)  \cdot2^{n}n!}{n!\cdot n!}\\
&  \ \ \ \ \ \ \ \ \ \ \ \ \ \ \ \ \ \ \ \ \left(  \text{since }\left(
2n\right)  !=\left(  1\cdot3\cdot5\cdot\cdots\cdot\left(  2n-1\right)
\right)  \cdot2^{n}n!\right) \\
&  =\dfrac{\left(  1\cdot3\cdot5\cdot\cdots\cdot\left(  2n-1\right)  \right)
\cdot2^{n}}{n!}=\dfrac{1\cdot3\cdot5\cdot\cdots\cdot\left(  2n-1\right)  }%
{n!}\cdot2^{n}.
\end{align*}
This proves Example \ref{exa.binom.2n-choose-n}.
\end{proof}

Entire books have been written about binomial coefficients and their
properties. See \cite{Spivey19} for a recent text (and \cite[Chapter 5]{GKP}
and \cite[Chapter 3]{detnotes} and \cite[\S 1.2.6]{Knuth-TAoCP1} and
\cite[Chapter 2]{Wildon19} for elementary introductions). Here are two more
basic facts that we will need (\cite[Theorem 1.3.8]{19fco} and
\cite[Proposition 1.3.6]{19fco}, respectively):

\begin{proposition}
[\emph{Pascal's identity}, aka \emph{recurrence of the binomial coefficients}%
]\label{prop.binom.rec}We have%
\begin{equation}
\dbinom{m}{n}=\dbinom{m-1}{n-1}+\dbinom{m-1}{n} \label{eq.binom.rec.m}%
\end{equation}
for any numbers $m$ and $n$.
\end{proposition}

\begin{proposition}
\label{prop.binom.0}Let $m,n\in\mathbb{N}$ satisfy $m<n$. Then, $\dbinom{m}%
{n}=0$.
\end{proposition}

Note that Proposition \ref{prop.binom.0} really requires $m\in\mathbb{N}$. For
example, $1.5<2$ but $\dbinom{1.5}{2}=0.375\neq0$.

Yet another useful property of the binomial coefficients is the following
(\cite[Theorem 1.3.11]{19fco}):

\begin{theorem}
[\emph{Symmetry of the binomial coefficients}]\label{thm.binom.sym}Let
$n\in\mathbb{N}$ and $k\in\mathbb{R}$. Then,%
\[
\dbinom{n}{k}=\dbinom{n}{n-k}.
\]

\end{theorem}

Note the $n\in\mathbb{N}$ requirement. Convince yourself that Theorem
\ref{thm.binom.sym} would fail for $n=-1$ and $k=0$.

\section{\label{chap.gf}Generating functions}

In this first chapter, we will discuss generating functions: first informally,
then on a rigorous footing. You may have seen generating functions already, as
their usefulness extends far beyond combinatorics; but they are so important
to this course that they are worth covering twice in case of doubt.

Rigorous introductions to generating functions (and formal power series in
general) can also be found in \cite[Chapter 7 (in the 1st edition)]{Loehr-BC},
in \cite[Chapter 1]{Henric74}, in \cite{Sambal22}, and (to some extent) in
\cite[Chapter 7]{19s}.\footnote{Bourbaki's \cite[\S IV.4]{Bourba03} contains
what might be the most rigorous and honest treatment of formal power series
available in the literature; however, it is not the most readable source, as
the notation is dense and heavily relies on other volumes by the same author.}
A quick overview is given in \cite{Niven69}, and many applications are found
in \cite{Wilf09}. There are furthermore numerous books that explore
enumerative combinatorics through the lens of generating functions
(\cite{GouJac83}, \cite{Wagner08}, \cite{Lando03} and others).


\subsection{\label{sec.gf.exas}Examples}

Let me first show what generating functions are good for. Then, starting in
the next section, I will explain how to rigorously define them. For now, I
will work informally; please suspend your disbelief until the next section.

The \textbf{idea} behind generating functions is easy: Any sequence $\left(
a_{0},a_{1},a_{2},\ldots\right)  $ of numbers gives rise to a
\textquotedblleft power series\textquotedblright\ $a_{0}+a_{1}x+a_{2}%
x^{2}+\cdots$, which is called the \emph{generating function} of this
sequence. This \textquotedblleft power series\textquotedblright\ is an
infinite sum (an \textquotedblleft infinite polynomial\textquotedblright\ in
an indeterminate $x$), so it is not immediately clear what it means and what
we are allowed to do with it; but before we answer such questions, let us
first play around with these power series and hope for the best. The following
four examples show how they can be useful.

\subsubsection{Example 1: The Fibonacci sequence}

\textbf{Example 1.} The \emph{Fibonacci sequence} is the sequence $\left(
f_{0},f_{1},f_{2},\ldots\right)  $ of integers defined recursively by%
\[
f_{0}=0,\ \ \ \ \ \ \ \ \ \ f_{1}=1,\ \ \ \ \ \ \ \ \ \ f_{n}=f_{n-1}%
+f_{n-2}\text{ for each }n\geq2.
\]
Its entries are known as the \emph{Fibonacci numbers}. Here are the first few
of them:%
\[%
\begin{tabular}
[c]{|c||c|c|c|c|c|c|c|c|c|c|c|c|}\hline
$n$ & $0$ & $1$ & $2$ & $3$ & $4$ & $5$ & $6$ & $7$ & $8$ & $9$ & $10$ &
$11$\\\hline
$f_{n}$ & $0$ & $1$ & $1$ & $2$ & $3$ & $5$ & $8$ & $13$ & $21$ & $34$ & $55$
& $89$\\\hline
\end{tabular}
\ \ \
\]


Let us see what we can learn about this sequence by considering its generating
function%
\begin{align*}
F\left(  x\right)   &  :=f_{0}+f_{1}x+f_{2}x^{2}+f_{3}x^{3}+\cdots\\
&  =0+1x+1x^{2}+2x^{3}+3x^{4}+5x^{5}+\cdots.
\end{align*}


We have%
\begin{align*}
F\left(  x\right)   &  =f_{0}+f_{1}x+f_{2}x^{2}+f_{3}x^{3}+f_{4}x^{4}+\cdots\\
&  =\underbrace{0+1x}_{=x}+\left(  f_{1}+f_{0}\right)  x^{2}+\left(
f_{2}+f_{1}\right)  x^{3}+\left(  f_{3}+f_{2}\right)  x^{4}+\cdots\\
&  \ \ \ \ \ \ \ \ \ \ \ \ \ \ \ \ \ \ \ \ \left(  \text{since }f_{0}=0\text{
and }f_{1}=1\text{ and }f_{n}=f_{n-1}+f_{n-2}\text{ for each }n\geq2\right) \\
&  =x+\underbrace{\left(  f_{1}+f_{0}\right)  x^{2}+\left(  f_{2}%
+f_{1}\right)  x^{3}+\left(  f_{3}+f_{2}\right)  x^{4}+\cdots}%
_{\substack{=\left(  f_{1}x^{2}+f_{2}x^{3}+f_{3}x^{4}+\cdots\right)  +\left(
f_{0}x^{2}+f_{1}x^{3}+f_{2}x^{4}+\cdots\right)  \\\text{(here we are hoping
that this manipulation of}\\\text{infinite sums is indeed legitimate)}}}\\
&  =x+\underbrace{\left(  f_{1}x^{2}+f_{2}x^{3}+f_{3}x^{4}+\cdots\right)
}_{\substack{=x\left(  f_{1}x+f_{2}x^{2}+f_{3}x^{3}+\cdots\right)  \\=x\left(
F\left(  x\right)  -f_{0}\right)  =xF\left(  x\right)  \\\text{(since }%
f_{0}=0\text{)}}}+\underbrace{\left(  f_{0}x^{2}+f_{1}x^{3}+f_{2}x^{4}%
+\cdots\right)  }_{\substack{=x^{2}\left(  f_{0}+f_{1}x+f_{2}x^{2}+f_{3}%
x^{3}+\cdots\right)  \\=x^{2}F\left(  x\right)  }}\\
&  =x+xF\left(  x\right)  +x^{2}F\left(  x\right)  =x+\left(  x+x^{2}\right)
F\left(  x\right)  .
\end{align*}
Solving this equation for $F\left(  x\right)  $ (assuming that we are allowed
to divide by $1-x-x^{2}$), we get%
\begin{align}
F\left(  x\right)   &  =\dfrac{x}{1-x-x^{2}}\label{eq.sec.gf.exas.1.Fx=1}\\
&  =\dfrac{x}{\left(  1-\phi_{+}x\right)  \left(  1-\phi_{-}x\right)
},\nonumber
\end{align}
where $\phi_{+}=\dfrac{1+\sqrt{5}}{2}$ and $\phi_{-}=\dfrac{1-\sqrt{5}}{2}$
are the two roots of the quadratic polynomial $1-x-x^{2}$ (note that $\phi
_{+}$ and $\phi_{-}$ are sometimes known as the \textquotedblleft golden
ratios\textquotedblright; we have $\phi_{+}\approx1.618$ and $\phi_{-}%
\approx-0.618$). Hence,%
\begin{align}
F\left(  x\right)   &  =\dfrac{x}{\left(  1-\phi_{+}x\right)  \left(
1-\phi_{-}x\right)  }\nonumber\\
&  =\dfrac{1}{\sqrt{5}}\cdot\dfrac{1}{1-\phi_{+}x}-\dfrac{1}{\sqrt{5}}%
\cdot\dfrac{1}{1-\phi_{-}x} \label{eq.sec.gf.exas.1.Fx=2}%
\end{align}
(by partial fraction decomposition).

Now, what are the coefficients of the power series $\dfrac{1}{1-\alpha x}$ for
an $\alpha\in\mathbb{C}$ ? Let me first answer this question for $\alpha=1$.
Namely, I claim that%
\begin{equation}
\dfrac{1}{1-x}=1+x+x^{2}+x^{3}+\cdots. \label{eq.sec.gf.exas.1.1/1-x}%
\end{equation}
Indeed, this follows by observing that
\begin{align*}
&  \left(  1-x\right)  \left(  1+x+x^{2}+x^{3}+\cdots\right) \\
&  =\left(  1+x+x^{2}+x^{3}+\cdots\right)  -x\left(  1+x+x^{2}+x^{3}%
+\cdots\right) \\
&  =\left(  1+x+x^{2}+x^{3}+\cdots\right)  -\left(  x+x^{2}+x^{3}+x^{4}%
+\cdots\right) \\
&  =1
\end{align*}
(again, we are hoping that these manipulations of infinite sums are allowed).
Note that the equality (\ref{eq.sec.gf.exas.1.1/1-x}) is a version of the
\emph{geometric series formula} familiar from real analysis. Now, for any
$\alpha\in\mathbb{C}$, we can substitute $\alpha x$ for $x$ in the equality
(\ref{eq.sec.gf.exas.1.1/1-x}), and thus obtain%
\begin{align}
\dfrac{1}{1-\alpha x}  &  =1+\alpha x+\left(  \alpha x\right)  ^{2}+\left(
\alpha x\right)  ^{3}+\cdots\nonumber\\
&  =1+\alpha x+\alpha^{2}x^{2}+\alpha^{3}x^{3}+\cdots.
\label{eq.sec.gf.exas.1.1/1-ax}%
\end{align}
Hence, our above formula (\ref{eq.sec.gf.exas.1.Fx=2}) becomes%
\begin{align*}
F\left(  x\right)   &  =\dfrac{1}{\sqrt{5}}\cdot\dfrac{1}{1-\phi_{+}x}%
-\dfrac{1}{\sqrt{5}}\cdot\dfrac{1}{1-\phi_{-}x}\\
&  =\dfrac{1}{\sqrt{5}}\cdot\left(  1+\phi_{+}x+\phi_{+}^{2}x^{2}+\phi_{+}%
^{3}x^{3}+\cdots\right)  -\dfrac{1}{\sqrt{5}}\cdot\left(  1+\phi_{-}x+\phi
_{-}^{2}x^{2}+\phi_{-}^{3}x^{3}+\cdots\right) \\
&  \ \ \ \ \ \ \ \ \ \ \ \ \ \ \ \ \ \ \ \ \left(  \text{by
(\ref{eq.sec.gf.exas.1.1/1-ax}), applied to }\alpha=\phi_{+}\text{ and again
to }\alpha=\phi_{-}\right) \\
&  =\dfrac{1}{\sqrt{5}}\sum_{k\geq0}\phi_{+}^{k}x^{k}-\dfrac{1}{\sqrt{5}}%
\sum_{k\geq0}\phi_{-}^{k}x^{k}\\
&  =\sum_{k\geq0}\left(  \dfrac{1}{\sqrt{5}}\cdot\phi_{+}^{k}-\dfrac{1}%
{\sqrt{5}}\cdot\phi_{-}^{k}\right)  x^{k}.
\end{align*}
Now, for any given $n\in\mathbb{N}$, the coefficient of $x^{n}$ in the power
series on the left hand side of this equality is $f_{n}$ (since $F\left(
x\right)  =f_{0}+f_{1}x+f_{2}x^{2}+f_{3}x^{3}+\cdots$), whereas the
coefficient of $x^{n}$ on the right hand side is clearly $\dfrac{1}{\sqrt{5}%
}\cdot\phi_{+}^{n}-\dfrac{1}{\sqrt{5}}\cdot\phi_{-}^{n}$. Thus, comparing
coefficients before $x^{n}$, we obtain%
\begin{align*}
f_{n}  &  =\dfrac{1}{\sqrt{5}}\cdot\phi_{+}^{n}-\dfrac{1}{\sqrt{5}}\cdot
\phi_{-}^{n}\\
&  \ \ \ \ \ \ \ \ \ \ \ \ \ \ \ \ \ \ \ \ \left(
\begin{array}
[c]{c}%
\text{assuming that \textquotedblleft comparing coefficients\textquotedblright%
\ is allowed, i.e.,}\\
\text{that equal power series really have equal coefficients}%
\end{array}
\right) \\
&  =\dfrac{1}{\sqrt{5}}\cdot\left(  \dfrac{1+\sqrt{5}}{2}\right)  ^{n}%
-\dfrac{1}{\sqrt{5}}\cdot\left(  \dfrac{1-\sqrt{5}}{2}\right)  ^{n}%
\end{align*}
for any $n\in\mathbb{N}$. This formula is known as \emph{Binet's formula}. It
has many consequences; for example, it implies easily that $\lim
\limits_{n\rightarrow\infty}\dfrac{f_{n+1}}{f_{n}}=\phi_{+}=\dfrac{1+\sqrt{5}%
}{2}\approx1.618\ldots$. Thus, $f_{n}\sim\phi_{+}^{n}$ in the asymptotical sense.

\subsubsection{Example 2: Dyck words and Catalan numbers}

Before the next example, let us address a warmup question: What is the number
of $2n$-tuples that contain $n$ entries equal to $0$ and $n$ entries equal to
$1$ (where $n$ is a given nonnegative integer)?

(For example, for $n=2$, these $2n$-tuples are $\left(  1,1,0,0\right)  $,
$\left(  1,0,1,0\right)  $, $\left(  1,0,0,1\right)  $, $\left(
0,1,1,0\right)  $, $\left(  0,1,0,1\right)  $, $\left(  0,0,1,1\right)  $, so
there are $6$ of them.)

\textbf{Answer:} The number is $\dbinom{2n}{n}$, since choosing a $2n$-tuple
that contains $n$ entries equal to $0$ and $n$ entries equal to $1$ is
tantamount to choosing an $n$-element subset of $\left\{  1,2,\ldots
,2n\right\}  $ (and we know that there are $\dbinom{2n}{n}$ ways to choose the latter).

\bigskip

\textbf{Example 2.} A \emph{Dyck word} of length $2n$ (where $n\in\mathbb{N}$)
is a $2n$-tuple that contains $n$ entries equal to $0$ and $n$ entries equal
to $1$, and has the additional property that for each $k$, we have
\begin{align}
&  \left(  \text{\# of }0\text{'s among its first }k\text{ entries}\right)
\nonumber\\
&  \leq\left(  \text{\# of }1\text{'s among its first }k\text{ entries}%
\right)  . \label{eq.sec.gf.exas.2.dyckword.leq}%
\end{align}
(The symbol \textquotedblleft\#\textquotedblright\ means \textquotedblleft
number\textquotedblright.)

Some examples: The tuples
\[
\left(  1,0,1,0\right)  ,\ \ \left(  1,1,0,0\right)  ,\ \ \left(
1,1,0,1,0,0\right)  ,\ \ \left(  {}\right)  ,\ \ \left(  1,0\right)
\]
are Dyck words. The tuples
\[
\left(  0,1,1,0\right)  ,\ \ \left(  1,0,0,1\right)  ,\ \ \left(
1,1,0\right)  ,\ \ \left(  1\right)  ,\ \ \left(  1,1,1,0\right)
\]
are not Dyck words.

A \emph{Dyck path} of length $2n$ is a path from the point $\left(
0,0\right)  $ to the point $\left(  2n,0\right)  $ in the Cartesian plane that
moves only using \textquotedblleft NE-steps\textquotedblright\ (i.e., steps of
the form $\left(  x,y\right)  \rightarrow\left(  x+1,y+1\right)  $) and
\textquotedblleft SE-steps\textquotedblright\ (i.e., steps of the form
$\left(  x,y\right)  \rightarrow\left(  x+1,y-1\right)  $) and never falls
below the x-axis (i.e., does not contain any point $\left(  x,y\right)  $ with
$y<0$).

Examples: For $n=2$, the Dyck paths from $\left(  0,0\right)  $ to $\left(
2n,0\right)  $ are%
\[%
\begin{tabular}
[c]{ccc}%
% [tikz figure removed]
%BeginExpansion
% [tikz figure removed]%
%EndExpansion
& $\ \ \ \ \ \ \ \ \ \ \ \ \ \ \ \ \ \ \ \ $ &
% [tikz figure removed]
%BeginExpansion
% [tikz figure removed]%
%EndExpansion
\end{tabular}
\]


A Dyck path can be viewed as the \textquotedblleft skyline\textquotedblright%
\ of a \textquotedblleft mountain range\textquotedblright. For example:%
\[%
\begin{tabular}
[c]{|c|c|}\hline
Dyck path & \textquotedblleft mountain range\textquotedblright\\\hline
\ \
% [tikz figure removed]
%BeginExpansion
% [tikz figure removed]%
%EndExpansion
\ \  & \ \
% [tikz figure removed]
%BeginExpansion
% [tikz figure removed]%
%EndExpansion
\ \ \\\hline
\end{tabular}
\]


The names \textquotedblleft NE-steps\textquotedblright\ and \textquotedblleft
SE-steps\textquotedblright\ in the definition of a Dyck path refer to compass
directions: If we treat the Cartesian plane as a map with the x-axis directed
eastwards and the y-axis directed northwards, then an NE-step moves to the
northeast, and an SE-step moves to the southeast.

Note that any NE-step and any SE-step increases the x-coordinate by $1$ (that
is, the step goes from a point with x-coordinate $k$ to a point with
x-coordinate $k+1$). Thus, any Dyck path from $\left(  0,0\right)  $ to
$\left(  2n,0\right)  $ has precisely $2n$ steps. Of these $2n$ steps, exactly
$n$ are NE-steps while the remaining $n$ are SE-steps (because any NE-step
increases the y-coordinate by $1$, while any SE-step decreases the
y-coordinate by $1$). Since a Dyck path must never fall below the x-axis, we
see that the number of SE-steps up to any given point can never be larger than
the number of NE-steps up to this point. But this is exactly the condition
(\ref{eq.sec.gf.exas.2.dyckword.leq}) from the definition of a Dyck word,
except that we are talking about NE-steps and SE-steps instead of $1$'s and
$0$'s. Thus, there is a simple bijection between Dyck words of length $2n$ and
Dyck paths from $\left(  0,0\right)  $ to $\left(  2n,0\right)  $:

\begin{itemize}
\item send each $1$ in the Dyck word to a NE-step in the Dyck path;

\item send each $0$ in the Dyck word to a SE-step in the Dyck path.
\end{itemize}

So the \# of Dyck words (of length $2n$) equals the \# of Dyck paths (from
$\left(  0,0\right)  $ to $\left(  2n,0\right)  $). But what is this number?

Example: For $n=3$, this number is $5$. Indeed, here are all Dyck paths from
$\left(  0,0\right)  $ to $\left(  6,0\right)  $, and their corresponding Dyck
words:%
\[%
\begin{tabular}
[c]{|c|c|}\hline
Dyck path & Dyck word\\\hline
\ \
% [tikz figure removed]
%BeginExpansion
% [tikz figure removed]%
%EndExpansion
\ \  & $\left(  1,1,0,0,1,0\right)  $\\\hline
\ \
% [tikz figure removed]
%BeginExpansion
% [tikz figure removed]%
%EndExpansion
\ \  & $\left(  1,1,1,0,0,0\right)  $\\\hline
\ \
% [tikz figure removed]
%BeginExpansion
% [tikz figure removed]%
%EndExpansion
\ \  & $\left(  1,0,1,0,1,0\right)  $\\\hline
\ \
% [tikz figure removed]
%BeginExpansion
% [tikz figure removed]%
%EndExpansion
\ \  & $\left(  1,0,1,1,0,0\right)  $\\\hline
\ \
% [tikz figure removed]
%BeginExpansion
% [tikz figure removed]%
%EndExpansion
\ \  & $\left(  1,1,0,1,0,0\right)  $\\\hline
\end{tabular}
\]
(We will soon stop writing the commas and parentheses when writing down words.
For example, the word $\left(  1,1,0,0,1,0\right)  $ will just become $110010$.)

Back to the general question.

For each $n\in\mathbb{N}$, we define
\begin{align*}
c_{n}:=  &  \ \left(  \text{\# of Dyck paths from }\left(  0,0\right)  \text{
to }\left(  2n,0\right)  \right) \\
=  &  \ \left(  \text{\# of Dyck words of length }2n\right)
\ \ \ \ \ \ \ \ \ \ \left(  \text{as we have seen above}\right)  .
\end{align*}
Then, $c_{0}=1$ (since the only Dyck path from $\left(  0,0\right)  $ to
$\left(  0,0\right)  $ is the trivial path) and $c_{1}=1$ and $c_{2}=2$ and
$c_{3}=5$ and $c_{4}=14$ and so on. These numbers $c_{n}$ are known as the
\emph{Catalan numbers}. Entire books have been written about them, such as
\cite{Stanle15}.

Let us first find a recurrence relation for $c_{n}$. The argument below is
best understood by following an example; namely, consider the following Dyck
path from $\left(  0,0\right)  $ to $\left(  16,0\right)  $ (so the
corresponding $n$ is $8$):%
\[%
% [tikz figure removed]
%BeginExpansion
% [tikz figure removed]%
%EndExpansion
\]


Fix a positive integer $n$. If $d$ is a Dyck path from $\left(  0,0\right)  $
to $\left(  2n,0\right)  $, then the \emph{first return} of $d$ (this is short
for \textquotedblleft first return of $d$ to the x-axis\textquotedblright)
shall mean the first point on $d$ that lies on the x-axis but is not the
origin (i.e., that has the form $\left(  i,0\right)  $ for some integer
$i>0$). For instance, in the example that we just gave, the first return is
the point $\left(  6,0\right)  $. If $d$ is a Dyck path from $\left(
0,0\right)  $ to $\left(  2n,0\right)  $, and if $\left(  i,0\right)  $ is its
first return, then $i$ is even\footnote{\textit{Proof.} The number of NE-steps
before the first return must equal the number of SE-steps before the first
return (because these steps have altogether taken us from the origin to a
point on the x-axis, and thus must have increased and decreased the
y-coordinate an equal number of times). This shows that the total number of
steps before the first return is even. In other words, $i$ is even (because
the total number of steps before the first return is $i$).}, and therefore we
have $i=2k$ for some $k\in\left\{  1,2,\ldots,n\right\}  $. Hence, for any
Dyck path from $\left(  0,0\right)  $ to $\left(  2n,0\right)  $, the first
return is a point of the form $\left(  2k,0\right)  $ for some $k\in\left\{
1,2,\ldots,n\right\}  $. Thus,%
\begin{align*}
&  \left(  \text{\# of Dyck paths from }\left(  0,0\right)  \text{ to }\left(
2n,0\right)  \right)  \\
&  =\sum_{k=1}^{n}\left(  \text{\# of Dyck paths from }\left(  0,0\right)
\text{ to }\left(  2n,0\right)  \text{ whose first return is }\left(
2k,0\right)  \right)  .
\end{align*}


Now, let us fix some $k\in\left\{  1,2,\ldots,n\right\}  $. We shall compute
the \# of Dyck paths from $\left(  0,0\right)  $ to $\left(  2n,0\right)  $
whose first return is $\left(  2k,0\right)  $. Any such Dyck path has a
natural \textquotedblleft two-part\textquotedblright\ structure: Its first
$2k$ steps form a path from $\left(  0,0\right)  $ to $\left(  2k,0\right)  $,
while its last (i.e., remaining) $2\left(  n-k\right)  $ steps form a path
from $\left(  2k,0\right)  $ to $\left(  2n,0\right)  $. Thus, in order to
construct such a path, we

\begin{itemize}
\item first choose its first $2k$ steps: They have to form a Dyck path from
$\left(  0,0\right)  $ to $\left(  2k,0\right)  $ that never returns to the
x-axis until $\left(  2k,0\right)  $. Hence, they begin with a NE-step and end
with a SE-step (since any other steps here would cause the path to fall below
the x-axis). Between these two steps, the remaining $2k-2=2\left(  k-1\right)
$ steps form a path that not only never falls below the x-axis, but also never
touches it (since $\left(  2k,0\right)  $ is the first return of our Dyck
path, so that our Dyck path does not touch the x-axis between $\left(
0,0\right)  $ and $\left(  2k,0\right)  $). In other words, these $2\left(
k-1\right)  $ steps form a path from $\left(  1,1\right)  $ to $\left(
2k-1,1\right)  $ that never falls below the $y=1$ line (i.e., below the x-axis
shifted by $1$ upwards). This means that it is a Dyck path from $\left(
0,0\right)  $ to $\left(  2\left(  k-1\right)  ,0\right)  $ (shifted by
$\left(  1,1\right)  $). Thus, there are $c_{k-1}$ possibilities for this
path. Hence, there are $c_{k-1}$ choices for the first $2k$ steps of our Dyck path.

\item then choose its last $2\left(  n-k\right)  $ steps: They have to form a
path from $\left(  2k,0\right)  $ to $\left(  2\left(  n-k\right)  ,0\right)
$ that never falls below the x-axis (but is allowed to touch it any number of
times). Thus, they form a Dyck path from $\left(  0,0\right)  $ to $\left(
2\left(  n-k\right)  ,0\right)  $ (shifted by $\left(  2k,0\right)  $). So
there are $c_{n-k}$ choices for these last $2\left(  n-k\right)  $ steps.
\end{itemize}

Thus, there are $c_{k-1}c_{n-k}$ many options for such a Dyck path from
$\left(  0,0\right)  $ to $\left(  2n,0\right)  $ (since choosing the first
$2k$ steps and choosing the last $2\left(  n-k\right)  $ steps are independent).

Let me illustrate this reasoning on the Dyck path from $\left(  0,0\right)  $
to $\left(  16,0\right)  $ shown above. This Dyck path has first return
$\left(  6,0\right)  $; thus, the corresponding $k$ is $3$. Since this Dyck
path does not return to the x-axis before $\left(  2k,0\right)  =\left(
6,0\right)  $, its first $2k$ steps stay above (or on) the yellow trapezoid
shown here:%
\[%
% [tikz figure removed]
%BeginExpansion
% [tikz figure removed]%
%EndExpansion
\]
In particular, the first and the last of these $2k$ steps are uniquely
determined, while the steps between them form a diagonally shifted Dyck path
that is filled in green here:%
\[%
% [tikz figure removed]
%BeginExpansion
% [tikz figure removed]%
%EndExpansion
\]
Finally, the last $2\left(  n-k\right)  $ steps form a horizontally shifted
Dyck path that is filled in purple here:%
\[%
% [tikz figure removed]
%BeginExpansion
% [tikz figure removed]%
%EndExpansion
\]
Our above argument shows that there are $c_{k-1}$ choices for the green Dyck
path and $c_{n-k}$ choices for the purple Dyck path, therefore $c_{k-1}%
c_{n-k}$ options in total.

Forget that we fixed $k$. Our counting argument above shows that%
\begin{align}
&  \left(  \text{\# of Dyck paths from }\left(  0,0\right)  \text{ to }\left(
2n,0\right)  \text{ whose first return is }\left(  2k,0\right)  \right)
\nonumber\\
&  =c_{k-1}c_{n-k} \label{eq.sec.gf.exas.2.dyckword.knum}%
\end{align}
for each $k\in\left\{  1,2,\ldots,n\right\}  $. Now,
\begin{align*}
c_{n}  &  =\left(  \text{\# of Dyck paths from }\left(  0,0\right)  \text{ to
}\left(  2n,0\right)  \right) \\
&  =\sum_{k=1}^{n}\underbrace{\left(  \text{\# of Dyck paths from }\left(
0,0\right)  \text{ to }\left(  2n,0\right)  \text{ whose first return is
}\left(  2k,0\right)  \right)  }_{\substack{=c_{k-1}c_{n-k}\\\text{(by
(\ref{eq.sec.gf.exas.2.dyckword.knum}))}}}\\
&  =\sum_{k=1}^{n}c_{k-1}c_{n-k}=c_{0}c_{n-1}+c_{1}c_{n-2}+c_{2}c_{n-3}%
+\cdots+c_{n-1}c_{0}.
\end{align*}
This is a recurrence equation for $c_{n}$. Combining it with $c_{0}=1$, we can
use it to compute any value of $c_{n}$ recursively. Let us, however, try to
digest it using generating functions!

Let%
\[
C\left(  x\right)  :=\sum_{n\geq0}c_{n}x^{n}=c_{0}+c_{1}x+c_{2}x^{2}%
+c_{3}x^{3}+\cdots.
\]
Thus,%
\begin{align*}
C\left(  x\right)   &  =c_{0}+c_{1}x+c_{2}x^{2}+c_{3}x^{3}+\cdots\\
&  =1+\left(  c_{0}c_{0}\right)  x+\left(  c_{0}c_{1}+c_{1}c_{0}\right)
x^{2}+\left(  c_{0}c_{2}+c_{1}c_{1}+c_{2}c_{0}\right)  x^{3}+\cdots\\
&  \ \ \ \ \ \ \ \ \ \ \ \ \ \ \ \ \ \ \ \ \left(
\begin{array}
[c]{c}%
\text{since }c_{0}=1\\
\text{and }c_{n}=c_{0}c_{n-1}+c_{1}c_{n-2}+c_{2}c_{n-3}+\cdots+c_{n-1}c_{0}\\
\text{for each }n>0
\end{array}
\right) \\
&  =1+x\underbrace{\left(  \left(  c_{0}c_{0}\right)  +\left(  c_{0}%
c_{1}+c_{1}c_{0}\right)  x+\left(  c_{0}c_{2}+c_{1}c_{1}+c_{2}c_{0}\right)
x^{2}+\cdots\right)  }_{\substack{=\left(  c_{0}+c_{1}x+c_{2}x^{2}%
+\cdots\right)  ^{2}\\\text{(because if we multiply out }\left(  c_{0}%
+c_{1}x+c_{2}x^{2}+\cdots\right)  ^{2}\\\text{and collect like powers of
}x\text{, we obtain}\\\text{exactly }\left(  c_{0}c_{0}\right)  +\left(
c_{0}c_{1}+c_{1}c_{0}\right)  x+\left(  c_{0}c_{2}+c_{1}c_{1}+c_{2}%
c_{0}\right)  x^{2}+\cdots\text{)}}}\\
&  =1+x\left(  \underbrace{c_{0}+c_{1}x+c_{2}x^{2}+\cdots}_{=C\left(
x\right)  }\right)  ^{2}=1+x\left(  C\left(  x\right)  \right)  ^{2}.
\end{align*}
This is a quadratic equation in $C\left(  x\right)  $. Let us solve it by the
quadratic formula (assuming for now that this is allowed -- i.e., that the
quadratic formula really does apply to our \textquotedblleft power
series\textquotedblright, whatever they are). Thus, we get%
\begin{equation}
C\left(  x\right)  =\dfrac{1\pm\sqrt{1-4x}}{2x}. \label{eq.sec.gf.exas.2.-+}%
\end{equation}
The $\pm$ sign here cannot be a $+$ sign, because if it was a $+$, then the
power series on top of the fraction would not be divisible by $2x$ (as its
constant term would be $2$ and thus nonzero\footnote{If you find this
unconvincing, here is a cleaner way to argue this: Multiplying the equality
(\ref{eq.sec.gf.exas.2.-+}) by $2x$, we obtain $2xC\left(  x\right)
=1\pm\sqrt{1-4x}$. The left hand side of this equality has constant term $0$,
but the right hand side has constant term $1\pm1$ (here, we are making the
assumption that $\sqrt{1-4x}$ is a power series with constant term $1$; this
is plausible because $\sqrt{1-4\cdot0}=1$ and will also be justified further
below). Thus, $0=1\pm1$; this shows that the $\pm$ sign is a $-$ sign.}).
Thus, (\ref{eq.sec.gf.exas.2.-+}) becomes%
\begin{equation}
C\left(  x\right)  =\dfrac{1-\sqrt{1-4x}}{2x}=\dfrac{1}{2x}\left(  1-\left(
1-4x\right)  ^{1/2}\right)  . \label{eq.sec.gf.exas.2.-}%
\end{equation}


How do we find the coefficients of the power series $\left(  1-4x\right)
^{1/2}$ ?

For each $n\in\mathbb{N}$, the binomial formula yields%
\begin{equation}
\left(  1+x\right)  ^{n}=\sum_{k=0}^{n}\dbinom{n}{k}x^{k}=\sum_{k\geq0}%
\dbinom{n}{k}x^{k}. \label{eq.sec.gf.exas.2.(1+x)n}%
\end{equation}
(Here, we have replaced the $\sum_{k=0}^{n}$ sign by a $\sum_{k\geq0}$ sign,
thus extending the summation from all $k\in\left\{  0,1,\ldots,n\right\}  $ to
all $k\in\mathbb{N}$. This does not change the value of the sum, since all the
newly appearing addends are $0$, as you can easily check.)

Let us pretend that the formula (\ref{eq.sec.gf.exas.2.(1+x)n}) holds not only
for $n\in\mathbb{N}$, but also for $n=1/2$. That is, we have%
\begin{equation}
\left(  1+x\right)  ^{1/2}=\sum_{k\geq0}\dbinom{1/2}{k}x^{k}.
\label{eq.sec.gf.exas.2.(1+x)1/2}%
\end{equation}
Now, substitute $-4x$ for $x$ in this equality (here we are making the rather
plausible assumption that we can substitute $-4x$ for $x$ in a power series);
then, we get%
\begin{align*}
\left(  1-4x\right)  ^{1/2}  &  =\sum_{k\geq0}\dbinom{1/2}{k}\left(
-4x\right)  ^{k}=\sum_{k\geq0}\dbinom{1/2}{k}\left(  -4\right)  ^{k}x^{k}\\
&  =\underbrace{\dbinom{1/2}{0}}_{=1}\underbrace{\left(  -4\right)  ^{0}}%
_{=1}\underbrace{x^{0}}_{=1}+\sum\limits_{k\geq1}\dbinom{1/2}{k}\left(
-4\right)  ^{k}x^{k}\\
&  =1+\sum\limits_{k\geq1}\dbinom{1/2}{k}\left(  -4\right)  ^{k}x^{k}.
\end{align*}
Hence,%
\[
1-\left(  1-4x\right)  ^{1/2}=1-\left(  1+\sum\limits_{k\geq1}\dbinom{1/2}%
{k}\left(  -4\right)  ^{k}x^{k}\right)  =-\sum\limits_{k\geq1}\dbinom{1/2}%
{k}\left(  -4\right)  ^{k}x^{k}.
\]
Thus, (\ref{eq.sec.gf.exas.2.-}) becomes%
\begin{align*}
C\left(  x\right)   &  =\dfrac{1}{2x}\underbrace{\left(  1-\left(
1-4x\right)  ^{1/2}\right)  }_{=-\sum\limits_{k\geq1}\dbinom{1/2}{k}\left(
-4\right)  ^{k}x^{k}}=\dfrac{1}{2x}\left(  -\sum\limits_{k\geq1}\dbinom
{1/2}{k}\left(  -4\right)  ^{k}x^{k}\right) \\
&  =\sum\limits_{k\geq1}\dbinom{1/2}{k}\underbrace{\dfrac{-\left(  -4\right)
^{k}x^{k}}{2x}}_{=2\left(  -4\right)  ^{k-1}x^{k-1}}=\sum\limits_{k\geq
1}\dbinom{1/2}{k}2\left(  -4\right)  ^{k-1}x^{k-1}\\
&  =\sum\limits_{k\geq0}\dbinom{1/2}{k+1}2\left(  -4\right)  ^{k}x^{k}\\
&  \ \ \ \ \ \ \ \ \ \ \ \ \ \ \ \ \ \ \ \ \left(  \text{here, we have
substituted }k+1\text{ for }k\text{ in the sum}\right)  .
\end{align*}
Comparing coefficients before $x^{n}$ in this equality gives%
\begin{equation}
c_{n}=\dbinom{1/2}{n+1}2\left(  -4\right)  ^{n}.
\label{eq.sec.gf.exas.cn=frac}%
\end{equation}


This is an explicit formula for $c_{n}$ (and makes computation of $c_{n}$
pretty easy!), but it turns out that it can be simplified further. Indeed, the
definition of $\dbinom{1/2}{n+1}$ yields%
\begin{align*}
\dbinom{1/2}{n+1}  &  =\dfrac{\left(  1/2\right)  \left(  1/2-1\right)
\left(  1/2-2\right)  \cdots\left(  1/2-n\right)  }{\left(  n+1\right)  !}\\
&  =\dfrac{\vphantom{\dfrac{f}{f}}\dfrac{1}{2}\cdot\dfrac{-1}{2}\cdot
\dfrac{-3}{2}\cdot\dfrac{-5}{2}\cdot\cdots\cdot\dfrac{-\left(  2n-1\right)
}{2}}{\left(  n+1\right)  !}\\
&  =\dfrac{\left(  1\cdot\left(  -1\right)  \cdot\left(  -3\right)
\cdot\left(  -5\right)  \cdot\cdots\cdot\left(  -\left(  2n-1\right)  \right)
\right)  /2^{n+1}}{\left(  n+1\right)  !}\\
&  =\dfrac{\left(  \left(  -1\right)  \cdot\left(  -3\right)  \cdot\left(
-5\right)  \cdot\cdots\cdot\left(  -\left(  2n-1\right)  \right)  \right)
/2^{n+1}}{\left(  n+1\right)  !}\\
&  =\dfrac{\left(  -1\right)  ^{n}\left(  1\cdot3\cdot5\cdot\cdots\cdot\left(
2n-1\right)  \right)  /2^{n+1}}{\left(  n+1\right)  !}.
\end{align*}
Thus, (\ref{eq.sec.gf.exas.cn=frac}) rewrites as
\begin{align*}
c_{n}  &  =\dfrac{\left(  -1\right)  ^{n}\left(  1\cdot3\cdot5\cdot\cdots
\cdot\left(  2n-1\right)  \right)  /2^{n+1}}{\left(  n+1\right)  !}%
\cdot2\left(  -4\right)  ^{n}\\
&  =\underbrace{\dfrac{1\cdot3\cdot5\cdot\cdots\cdot\left(  2n-1\right)
}{\left(  n+1\right)  !}}_{\substack{=\dfrac{1}{n+1}\cdot\dfrac{1\cdot
3\cdot5\cdot\cdots\cdot\left(  2n-1\right)  }{n!}\\\text{(since }\left(
n+1\right)  !=\left(  n+1\right)  \cdot n!\text{)}}}\cdot\,\underbrace{\dfrac
{\left(  -1\right)  ^{n}\cdot2\left(  -4\right)  ^{n}}{2^{n+1}}}_{=2^{n}}\\
&  =\dfrac{1}{n+1}\cdot\underbrace{\dfrac{1\cdot3\cdot5\cdot\cdots\cdot\left(
2n-1\right)  }{n!}\cdot2^{n}}_{\substack{=\dbinom{2n}{n}\\\text{(by Example
\ref{exa.binom.2n-choose-n})}}}=\dfrac{1}{n+1}\dbinom{2n}{n}.
\end{align*}
Hence, we have shown that
\begin{equation}
c_{n}=\dfrac{1}{n+1}\dbinom{2n}{n}. \label{eq.sec.gf.exas.cn=1/(n+1)}%
\end{equation}
Moreover, we can rewrite this further as%
\begin{equation}
c_{n}=\dbinom{2n}{n}-\dbinom{2n}{n-1} \label{eq.sec.gf.exas.cn=diff}%
\end{equation}
(since another binomial coefficient manipulation\footnote{See Exercise
\ref{exe.gf.cn-easy-manip} \textbf{(a)} for this.} yields $\dfrac{1}%
{n+1}\dbinom{2n}{n}=\dbinom{2n}{n}-\dbinom{2n}{n-1}$).

Here is the upshot: The \# of Dyck words of length $2n$ is $c_{n}=\dfrac
{1}{n+1}\dbinom{2n}{n}$. In other words, a $2n$-tuple that consists of $n$
entries equal to $0$ and $n$ entries equal to $1$ (chosen uniformly at random)
is a Dyck word with probability $\dfrac{1}{n+1}$.

(There are also combinatorial ways to prove this; see, e.g., \cite[\S 7.5,
discussion at the end of Example 4]{GKP} or \cite[\S 1.6]{Stanle15} or
\cite{Martin-Dyck} or \cite[Lecture 29, Theorem 5.6.7]{22fco} or \cite[Theorem
1.56]{Loehr-BC}\footnote{Note that the \textquotedblleft Dyck
paths\textquotedblright\ in \cite{Loehr-BC} differ from ours in that they use
N-steps (i.e., steps $\left(  i,j\right)  \mapsto\left(  i,j+1\right)  $) and
E-steps (i.e., steps $\left(  i,j\right)  \mapsto\left(  i+1,j\right)  $)
instead of NE-steps and SE-steps, and stay above the $x=y$ line instead of
above the x-axis. But this notion of Dyck paths is equivalent to ours, since a
clockwise rotation by $45^{\circ}$ followed by a $\sqrt{2}$-homothety
transforms it into ours.} or \cite[\S 8.5, proofs of Identity 244]%
{Spivey19}\footnote{Again, \cite{Spivey19} works not directly with Dyck paths,
but rather with paths that use E-steps (i.e., steps $\left(  i,j\right)
\mapsto\left(  i+1,j\right)  $) and N-steps (i.e., steps $\left(  i,j\right)
\mapsto\left(  i,j+1\right)  $) instead of NE-steps and SE-steps, and stay
below the $x=y$ line instead of above the x-axis. But this kind of Dyck paths
is equivalent to our Dyck paths, since a reflection across the $x=y$ line,
followed by a clockwise rotation by $45^{\circ}$ followed by a $\sqrt{2}%
$-homothety transforms it into ours.}.)

Here is a table of the first $12$ Catalan numbers $c_{n}$:%
\[%
\begin{tabular}
[c]{|c||c|c|c|c|c|c|c|c|c|c|c|c|}\hline
$n$ & $0$ & $1$ & $2$ & $3$ & $4$ & $5$ & $6$ & $7$ & $8$ & $9$ & $10$ &
$11$\\\hline
$c_{n}$ & $1$ & $1$ & $2$ & $5$ & $14$ & $42$ & $132$ & $429$ & $1430$ &
$4862$ & $16\,796$ & $58\,786$\\\hline
\end{tabular}
\ \ \ \ \ .
\]


\subsubsection{Example 3: The Vandermonde convolution}

\textbf{Example 3:} The \emph{Vandermonde convolution identity} (also known as
the \emph{Chu--Vandermonde identity}) says that%
\[
\dbinom{a+b}{n}=\sum_{k=0}^{n}\dbinom{a}{k}\dbinom{b}{n-k}%
\ \ \ \ \ \ \ \ \ \ \text{for any numbers }a,b\text{ and any }n\in\mathbb{N}%
\]
(where \textquotedblleft numbers\textquotedblright\ can mean, e.g.,
\textquotedblleft complex numbers\textquotedblright).

Let us prove this using generating functions. For now, we shall only prove
this for $a,b\in\mathbb{N}$; later I will explain why it also holds for
arbitrary (rational, real or complex) numbers $a,b$ as well.

Indeed, fix $a,b\in\mathbb{N}$. Recall (from (\ref{eq.sec.gf.exas.2.(1+x)n}))
that
\[
\left(  1+x\right)  ^{n}=\sum_{k\geq0}\dbinom{n}{k}x^{k}%
\]
for each $n\in\mathbb{N}$. Hence,%
\begin{align*}
\left(  1+x\right)  ^{a}  &  =\sum_{k\geq0}\dbinom{a}{k}x^{k}%
\ \ \ \ \ \ \ \ \ \ \text{and}\\
\left(  1+x\right)  ^{b}  &  =\sum_{k\geq0}\dbinom{b}{k}x^{k}%
\ \ \ \ \ \ \ \ \ \ \text{and}\\
\left(  1+x\right)  ^{a+b}  &  =\sum_{k\geq0}\dbinom{a+b}{k}x^{k}.
\end{align*}
Thus,
\begin{align*}
\underbrace{\left(  1+x\right)  ^{a}}_{=\sum_{k\geq0}\dbinom{a}{k}x^{k}%
}\ \ \underbrace{\left(  1+x\right)  ^{b}}_{=\sum_{k\geq0}\dbinom{b}{k}x^{k}}
&  =\left(  \sum_{k\geq0}\dbinom{a}{k}x^{k}\right)  \left(  \sum_{k\geq
0}\dbinom{b}{k}x^{k}\right) \\
&  =\left(  \sum_{k\geq0}\dbinom{a}{k}x^{k}\right)  \left(  \sum_{\ell\geq
0}\dbinom{b}{\ell}x^{\ell}\right) \\
&  =\sum_{k\geq0}\ \ \sum_{\ell\geq0}\dbinom{a}{k}x^{k}\dbinom{b}{\ell}%
x^{\ell}\\
&  =\sum_{k\geq0}\ \ \sum_{\ell\geq0}\dbinom{a}{k}\dbinom{b}{\ell}x^{k+\ell}\\
&  =\sum_{n\geq0}\left(  \sum_{k=0}^{n}\dbinom{a}{k}\dbinom{b}{n-k}\right)
x^{n}%
\end{align*}
(here, we have merged addends in which $x$ appears in the same power). Hence,%
\begin{align*}
\sum_{n\geq0}\left(  \sum_{k=0}^{n}\dbinom{a}{k}\dbinom{b}{n-k}\right)  x^{n}
&  =\left(  1+x\right)  ^{a}\left(  1+x\right)  ^{b}=\left(  1+x\right)
^{a+b}=\sum_{k\geq0}\dbinom{a+b}{k}x^{k}\\
&  =\sum_{n\geq0}\dbinom{a+b}{n}x^{n}.
\end{align*}
Comparing coefficients in this equality, we obtain%
\[
\sum_{k=0}^{n}\dbinom{a}{k}\dbinom{b}{n-k}=\dbinom{a+b}{n}%
\ \ \ \ \ \ \ \ \ \ \text{for each }n\in\mathbb{N}.
\]
This completes the proof of the Vandermonde convolution identity for
$a,b\in\mathbb{N}$.

\subsubsection{Example 4: Solving a recurrence}

\textbf{Example 4.} The following example is from \cite[\S 1.2]{Wilf04}.
Define a sequence $\left(  a_{0},a_{1},a_{2},\ldots\right)  $ of numbers
recursively by%
\[
a_{0}=1,\ \ \ \ \ \ \ \ \ \ a_{n+1}=2a_{n}+n\text{ for all }n\geq0.
\]
Thus, the first entries of this sequence are $1,2,5,12,27,58,121,\ldots$. This
sequence appears in \href{https://oeis.org/}{the OEIS (= Online Encyclopedia
of Integer Sequences)} as \href{https://oeis.org/A000325}{A000325}, with index shifted.

Can we find an explicit formula for $a_{n}$ (without looking it up in the OEIS)?

Again, generating functions are helpful. Set%
\[
A\left(  x\right)  =a_{0}+a_{1}x+a_{2}x^{2}+a_{3}x^{3}+\cdots.
\]
Then,%
\begin{align}
A\left(  x\right)   &  =a_{0}+a_{1}x+a_{2}x^{2}+a_{3}x^{3}+\cdots\nonumber\\
&  =1+\left(  2a_{0}+0\right)  x+\left(  2a_{1}+1\right)  x^{2}+\left(
2a_{2}+2\right)  x^{3}+\cdots\nonumber\\
&  \ \ \ \ \ \ \ \ \ \ \ \ \ \ \ \ \ \ \ \ \left(  \text{since }a_{0}=1\text{
and }a_{n+1}=2a_{n}+n\text{ for all }n\geq0\right) \nonumber\\
&  =1+2\underbrace{\left(  a_{0}x+a_{1}x^{2}+a_{2}x^{3}+\cdots\right)
}_{=xA\left(  x\right)  }+\underbrace{\left(  0x+1x^{2}+2x^{3}+\cdots\right)
}_{=x\left(  0+1x+2x^{2}+3x^{3}+\cdots\right)  }\nonumber\\
&  =1+2xA\left(  x\right)  +x\left(  0+1x+2x^{2}+3x^{3}+\cdots\right)  .
\label{eq.sec.gf.exas.4.Ax=1}%
\end{align}


Thus, it would clearly be helpful to find a simple expression for
$0+1x+2x^{2}+3x^{3}+\cdots$. Here are two ways to do so: \medskip

\textit{First way:} We assume that our power series (whatever they actually
are) can be differentiated (as if they were functions). We furthermore assume
that these derivatives satisfy the same basic rules (sum rule, product rule,
quotient rule, chain rule) as the derivatives in real analysis. (Again, these
assumptions shall be justified later on.)

Denoting the derivative of a power series $f$ by $f^{\prime}$, we then have%
\[
\left(  1+x+x^{2}+x^{3}+\cdots\right)  ^{\prime}=1+2x+3x^{2}+4x^{3}+\cdots.
\]
Hence,%
\[
1+2x+3x^{2}+4x^{3}+\cdots=\left(  1+x+x^{2}+x^{3}+\cdots\right)  ^{\prime
}=\left(  \dfrac{1}{1-x}\right)  ^{\prime}%
\]
(since (\ref{eq.sec.gf.exas.1.1/1-x}) yields $1+x+x^{2}+x^{3}+\cdots=\dfrac
{1}{1-x}$). Using the quotient rule, we can easily find that $\left(
\dfrac{1}{1-x}\right)  ^{\prime}=\dfrac{1}{\left(  1-x\right)  ^{2}}$, so that%
\begin{equation}
1+2x+3x^{2}+4x^{3}+\cdots=\left(  \dfrac{1}{1-x}\right)  ^{\prime}=\dfrac
{1}{\left(  1-x\right)  ^{2}}. \label{eq.sec.gf.exas.4.1/(1-x)2}%
\end{equation}


The left hand side of this looks very similar to the power series
$0+1x+2x^{2}+3x^{3}+\cdots$ that we want to simplify. And indeed, we have the
following:%
\begin{align}
0+1x+2x^{2}+3x^{3}+\cdots &  =x\underbrace{\left(  1+2x+3x^{2}+4x^{3}%
+\cdots\right)  }_{=\dfrac{1}{\left(  1-x\right)  ^{2}}}=x\cdot\dfrac
{1}{\left(  1-x\right)  ^{2}}\nonumber\\
&  =\dfrac{x}{\left(  1-x\right)  ^{2}}. \label{eq.sec.gf.exas.4.x/(1-x)2}%
\end{align}
\medskip

\textit{Second way:} We rewrite $0+1x+2x^{2}+3x^{3}+\cdots$ as an infinite sum
of infinite sums:%
\begin{align}
&  \ 0+1x+2x^{2}+3x^{3}+\cdots\nonumber\\
&
\begin{array}
[c]{cccccccccc}%
= & x^{1} & + & x^{2} & + & x^{3} & + & x^{4} & + & \cdots\\
&  & + & x^{2} & + & x^{3} & + & x^{4} & + & \cdots\\
&  &  &  & + & x^{3} & + & x^{4} & + & \cdots\\
&  &  &  &  &  & + & x^{4} & + & \cdots\\
&  &  &  &  &  &  &  & \ddots & \ddots
\end{array}
\nonumber\\
&  =\sum_{k\geq1}\underbrace{\left(  x^{k}+x^{k+1}+x^{k+2}+\cdots\right)
}_{\substack{=x^{k}\cdot\left(  1+x+x^{2}+x^{3}+\cdots\right)  \\=x^{k}%
\cdot\dfrac{1}{1-x}\\\text{(by (\ref{eq.sec.gf.exas.1.1/1-x}))}}}=\sum
_{k\geq1}x^{k}\cdot\dfrac{1}{1-x}=\dfrac{1}{1-x}\cdot\underbrace{\sum_{k\geq
1}x^{k}}_{\substack{=x^{1}+x^{2}+x^{3}+\cdots\\=x\cdot\left(  1+x+x^{2}%
+x^{3}+\cdots\right)  \\=x\cdot\dfrac{1}{1-x}\\\text{(by
(\ref{eq.sec.gf.exas.1.1/1-x}))}}}\nonumber\\
&  =\dfrac{1}{1-x}\cdot x\cdot\dfrac{1}{1-x}=\dfrac{x}{\left(  1-x\right)
^{2}}. \label{eq.sec.gf.exas.4.x/(1-x)2alt}%
\end{align}
(We used some unstated assumptions here about infinite sums -- specifically,
we assumed that we can rearrange them without worrying about absolute
convergence or similar issues -- but we will later see that these assumptions
are well justified. Besides, we obtained the same result as by our first way,
which is reassuring.) \medskip

Having computed $0+1x+2x^{2}+3x^{3}+\cdots$, we can now simplify
(\ref{eq.sec.gf.exas.4.Ax=1}), obtaining%
\begin{align*}
A\left(  x\right)   &  =1+2xA\left(  x\right)  +x\underbrace{\left(
0+1x+2x^{2}+3x^{3}+\cdots\right)  }_{=\dfrac{x}{\left(  1-x\right)  ^{2}}}\\
&  =1+2xA\left(  x\right)  +x\cdot\dfrac{x}{\left(  1-x\right)  ^{2}}.
\end{align*}
This is a linear equation in $A\left(  x\right)  $. Solving it yields%
\begin{align*}
A\left(  x\right)   &  =\dfrac{1-2x+2x^{2}}{\left(  1-x\right)  ^{2}\left(
1-2x\right)  }\\
&  =\underbrace{\dfrac{2}{1-2x}}_{\substack{=2\sum_{k\geq0}2^{k}%
x^{k}\\\text{(by (\ref{eq.sec.gf.exas.1.1/1-ax}))}}}-\underbrace{\dfrac
{1}{\left(  1-x\right)  ^{2}}}_{\substack{=1+2x+3x^{2}+4x^{3}+\cdots
\\\text{(by (\ref{eq.sec.gf.exas.4.1/(1-x)2}), or alternatively}\\\text{by
dividing the}\\\text{equality (\ref{eq.sec.gf.exas.4.x/(1-x)2alt}) by
}x\text{)}}}\ \ \ \ \ \ \ \ \ \ \left(  \text{by partial fraction
decomposition}\right) \\
&  =\underbrace{2\sum_{k\geq0}2^{k}x^{k}}_{=\sum_{k\geq0}2^{k+1}x^{k}%
}-\underbrace{\left(  1+2x+3x^{2}+4x^{3}+\cdots\right)  }_{=\sum_{k\geq
0}\left(  k+1\right)  x^{k}}\\
&  =\sum_{k\geq0}2^{k+1}x^{k}-\sum_{k\geq0}\left(  k+1\right)  x^{k}%
=\sum_{k\geq0}\left(  2^{k+1}-\left(  k+1\right)  \right)  x^{k}.
\end{align*}
Comparing coefficients, we obtain%
\[
a_{n}=2^{n+1}-\left(  n+1\right)  \ \ \ \ \ \ \ \ \ \ \text{for each }%
n\in\mathbb{N}.
\]
This is also easy to prove directly (by induction on $n$).
